
A model achieving 95\% accuracy on ImageNet may fail 40\% of the time on ObjectNet, yet no quantitative metric predicts which benchmarks will exhibit such generalization gaps. Recht et al. (2019) demonstrated that top-performing ImageNet classifiers suffer 11-15\% accuracy drops when evaluated on ImageNet-V2, a carefully reproduced test set following the original data collection methodology. McCoy et al. (2019) showed that BERT achieves 84\% on MNLI but performs near chance on HANS when syntactic heuristics are isolated. These generalization gaps represent research effort directed at improvements that do not transfer to deployment conditions.

Prior work has focused on measuring gaps after they occur, not predicting them. Recht et al. quantified the ImageNet gap but offered no metric to identify saturated benchmarks a priori. ObjectNet and HANS exposed model failures but required constructing new test sets to reveal them.

This paper proposes the Difficulty-Normalized Saturation Index (DNSI), based on the observation that benchmark saturation has an information-theoretic signature. As benchmarks mature, the distribution of performance improvements shifts: early progress is diverse (many approaches, substantial gains), while late-stage progress is homogeneous (minor tweaks, incremental gains). This shift is measurable as entropy of improvement patterns. DNSI is defined as the ratio of observed improvement entropy to expected entropy based on task difficulty, separating true saturation from benchmark hardness.

The contributions are:

\begin{enumerate}
\item DNSI, an entropy-based saturation metric with difficulty normalization, computable from publicly available SOTA histories without requiring held-out test sets.
\end{enumerate}

\begin{enumerate}
\item Correlation analysis showing DNSI relates to generalization gaps (Pearson R = -0.950, p = 0.050, n = 4) across four benchmarks with published held-out evaluations.
\end{enumerate}

\begin{enumerate}
\item Preliminary evidence that pre-saturation DNSI values correlate with future generalization gaps ($R^2$ = 0.349), with consistent negative direction across vision (R = -0.972) and NLP (R = -0.684) domains.
\end{enumerate}
